
\section{约定与预备知识}
\label{sec:8-2}

本节将固定若干基本约定，并准备一些关于定向的公式。开始时，我们假定
所论空间都是光滑空间（即流形），映射都是浸没。随后再纳入仓西结构，
讨论带仓西结构空间以及仓西结构意义下弱浸没强光滑映射的纤维积。
本节用相应的小写字母 $x$ 或 $y$ 表示空间 $X$ 或 $Y$ 的维数。
从现在起，我们以 $*$ 表示一个不再具体指明的一般点。

\newpage
\begin{statement}{约定}{8.2.1}\label{conv:8-2-1}
\begin{enumerate}[label=(\arabic*),leftmargin=2.8em]
\item 设 $X$ 是带边界 $\partial X$ 的定向光滑流形。%
\liftIndexRangeStart{I299}我们规定
$\partial X$ 的定向，使
\begin{equation}
 T_*X\cong\mathbb R_{\mathrm{out}}\times T_*(\partial X)
 \tag{8.2.1.1}\label{eq:8-2-1-1}
\end{equation}
是定向向量空间的同构。这里
$\mathbb R_{\mathrm{out}}$ 是由外法向量定向的 $\mathbb R$。

为了讨论带仓西结构空间的情形，先准备一些记号（详见第 A.1 节）。
依定义 A1.1，$p\in X$ 周围的一个仓西邻域是五元组
$(V_p,E_p,\Gamma_p,\psi_p,s_p)$；记商空间 $V_p/\Gamma_p$ 为 $U_p$。
按照仓西结构定向的定义（见定义 A1.17 及 \textbf{[FuOn99II]}
定义 5.8），$\det E_p^*\otimes\det T_*V_p$ 带有定向，而有限群
$\Gamma_p$ 在 $E_p$ 和 $V_p$ 上的作用保持
$\det E_p^*\otimes\det T_*V_p$ 的定向。因此群作用不影响这里的
定向问题；在以下论证中，可以假定该作用平凡。换言之，我们考察
$\Gamma_p$ 平凡的情形，并把 $V_p$ 与 $U_p$ 识别。

本章把一个（定向）仓西邻域，也就是带定向仓西结构空间的局部坐标图，
记作
\[
 (s;E\longrightarrow U),
\]
其中 $E\to U$ 是向量丛，$s$ 是它的一个截面。换言之，$s$ 是 $U$ 上
向量丛 $E$ 的一个截面。当 $s$ 不横截于零截面时，
我们用多值扰动得到虚基本链。这里的任务是给每个零点集赋予典范定向，
故假定经过扰动后 $s$ 横截于零截面。于是有正合列
\begin{equation}
 0\longrightarrow T_*s^{-1}(0)\longrightarrow T_*U
 \xrightarrow{D_*s}E_*\longrightarrow0.
 \tag{8.2.2}\label{eq:8-2-2}
\end{equation}
这里 $D_*s\colon T_*U\to E_*$ 是借助一个联络取得的
$s\colon U\to E$ 在 $*$ 处的协变导数。只要 $*\in s^{-1}(0)$，
这一导数就不依赖于联络的选择。我们排列 $T_*U$ 的基，使 $E_*$ 的基
在前，$T_*s^{-1}(0)$ 的基在后，并由
\[
 E_*\times T_*s^{-1}(0)=T_*U
\]
定义 $s^{-1}(0)$ 的定向。（由于
$\det(E_*)^*\otimes\det T_*U$ 的定向已经给定，上式确实确定
$T_*s^{-1}(0)$ 的定向。）

设 $(s;E\to U)$ 是带边定向仓西结构空间的一个局部坐标图。回忆
$(s;E\to U)$ 的定向就是 $\det E^*\otimes\det TU$ 的定向。
必要时缩小 $U$，便可假定 $U$ 是带边定向流形，而 $E$ 是其上的
定向向量丛。我们把 $(s;E\to U)$ 的边界定向定义为
\[
 \partial(s;E\longrightarrow U)
 =(-1)^{\operatorname{rank}E}
 \bigl(s|_{\partial U};E|_{\partial U}\longrightarrow\partial U\bigr),
\]
其中 $\partial U$ 带有 $U$ 的边界定向。
\liftIndexRangeEnd{I299}

\item 设李群 $G$ 带有一个定向。当 $G$ 光滑、自由地右作用在定向流形
$X$ 上时，我们规定商空间 $X/G$\liftIndex{I300} 的定向，使
\begin{equation}
 T_*X\cong T_*(X/G)\times\operatorname{Lie}G
 \tag{8.2.1.2}\label{eq:8-2-1-2}
\end{equation}
是定向向量空间的同构。这里 $\operatorname{Lie}G$ 是 $G$ 的李代数。

若 $G$ 作用在带定向仓西结构的空间 $X$ 上，并且每一点的稳定子群都
有限，则商空间继承定向仓西结构。具体地，对一个 $G$-等变仓西邻域
$(s;E\to U)$，商空间 $U/G$ 如上定向。向量丛 $E$ 及其截面 $s$
下降为（轨形）丛 $E/G\to U/G$ 及其截面 $s$。于是
$(s;E/G\to U/G)$ 带有定向，因为对 $x\in U$ 及其等价类
$[x]\in U/G$，有
\[
\begin{aligned}
 \det E_x^*\otimes\det T_xU
 &\cong \det E_{[x]}^*\otimes
   \det\bigl(T_{[x]}(U/G)\oplus\operatorname{Lie}G\bigr)\\
 &\cong \det E_{[x]}^*\otimes\det T_{[x]}(U/G)
   \otimes\det\operatorname{Lie}G.
\end{aligned}
\]

\item \textup{（纤维积定向\liftIndex{I301}）}
设 $X_i$（$i=1,2$）和 $Y$ 是定向光滑流形，
$f_i\colon X_i\to Y$ 是两个浸没。把纤维积记作
\[
 X_1{}_{f_1}\!\times_{f_2}X_2
 \quad\text{或}\quad X_1\times_YX_2.
\]
取 $q\in f_1(X_1)\cap f_2(X_2)$，并设
$(u_1,\ldots,u_y)$ 是 $T_qY$ 的一组定向基。由横截性，对
$p_i\in f_i^{-1}(q)$，可以选取
\[
 (v_{X_i,1},\ldots,v_{X_i,y})\subset T_{p_i}X_i,
 \qquad
 (df_i)_{p_i}(v_{X_i,k})=u_k
 \quad(k=1,\ldots,y).
\]
还可以在 $\ker(df_i)_{p_i}$ 中选取基
$(\eta_{X_i,1},\ldots,\eta_{X_i,x_i-y})$，使
\[
 (\eta_{X_1,1},\ldots,\eta_{X_1,x_1-y},
   v_{X_1,1},\ldots,v_{X_1,y})
\]
与 $T_{p_1}X_1$ 的给定定向相容，其中
\[
 T_{p_1}X_1\cong
 \ker(df_1)_{p_1}\times\operatorname{Im}(df_1)_{p_1},
\]
而
\[
 (v_{X_2,1},\ldots,v_{X_2,y},
   \eta_{X_2,1},\ldots,\eta_{X_2,x_2-y})
\]
与 $T_{p_2}X_2$ 的给定定向相容，其中
\[
 T_{p_2}X_2\cong
 \operatorname{Im}(df_2)_{p_2}\times\ker(df_2)_{p_2}.
\]
我们规定纤维积 $X_1\times_YX_2$ 的定向，使
\[
 (\eta_{X_1,1},\ldots,\eta_{X_1,x_1-y},u_1,\ldots,u_y,
   \eta_{X_2,1},\ldots,\eta_{X_2,x_2-y})
\]
是 $T_{[p_1,p_2]}(X_1\times_YX_2)$ 的一组定向基。这里借助
$f_i$ 的局部截面 $s_i$，把 $T_qY$ 与
$T_{[p_1,p_2]}((s_1\times s_2)Y)$ 识别。

换言之，先把向量空间
$T_{[p_1,p_2]}(X_1\times_YX_2)$ 识别为
\[
 \ker(df_1)_{p_1}\times T_qY\times\ker(df_2)_{p_2},
\]
再规定其定向，使这一识别为定向向量空间的同构。这里
$\ker(df_i)_{p_i}$ 的定向由下列定向同构分别规定：
\begin{align}
 T_{p_1}X_1
 &\cong\ker(df_1)_{p_1}\times\operatorname{Im}(df_1)_{p_1},
 \tag{8.2.1.3}\label{eq:8-2-1-3}\\
 T_{p_2}X_2
 &\cong\operatorname{Im}(df_2)_{p_2}\times\ker(df_2)_{p_2}.
 \tag{8.2.1.4}\label{eq:8-2-1-4}
\end{align}
其中 $\operatorname{Im}(df_i)_{p_i}$ 与 $T_qY$ 作为定向向量空间同构。
当 $(df_i)_{p_i}$（$i=1,2$）为双射时，我们用 $\epsilon_i$ 给单点集
$\ker(df_i)_{p_i}$ 定向，并由此记录定向同构\liftIndex{I302}：
若 $(df_i)_{p_i}$ 是保定向同构，则
$\epsilon_i=+1$；若它是反定向同构，则 $\epsilon_i=-1$。

按照这种写法，今后讨论纤维积定向时，将把定向同构
\eqref{eq:8-2-1-3}--\eqref{eq:8-2-1-4} 简写为
\begin{align}
 X_1&=X_1^\circ\times Y,
 \tag{8.2.1.5}\label{eq:8-2-1-5}\\
 X_2&=Y\times{}^\circ X_2.
 \tag{8.2.1.6}\label{eq:8-2-1-6}
\end{align}
相应地，把纤维积 $X_1\times_YX_2$ 写成
$X_1^\circ\times Y\times{}^\circ X_2$。当然，这些等式是在切空间层次
成立的等式。

\item \textup{（仓西结构的纤维积）}
设 $(s_1;E_1\to U_1)$ 和 $(s_2;E_2\to U_2)$ 分别是
$X_1$ 和 $X_2$ 的仓西邻域。可以取乘积
\[
 (s_1\oplus s_2;E_1\oplus E_2\longrightarrow U_1\times U_2),
\]
它给出 $X_1\times X_2$ 上的仓西结构。注意，由这一邻域在
$(s_1\oplus s_2)^{-1}(0)$ 上定义的定向不同于
\[
 X_1\times X_2=s_1^{-1}(0)\times s_2^{-1}(0)
\]
的乘积定向；两者的差为
\[
\begin{aligned}
 X_1\times X_2
 &=(-1)^{\operatorname{rank}E_2\dim X_1}
   (s_1\oplus s_2)^{-1}(0)\\
 &=(-1)^{\operatorname{rank}E_2(\dim U_1-\operatorname{rank}E_1)}
   (s_1\oplus s_2)^{-1}(0).
\end{aligned}
\]

一般纤维积的情形如下定义。假定定义纤维积的映射都是弱浸没的
（见第 A1 节），并设
$f_1\colon U_1\to B$、$f_2\colon U_2\to B$ 是浸没。取纤维积
\[
 U_1\times_BU_2=U_1{}_{f_1}\!\times_{f_2}U_2,
\]
并按本约定 (3) 赋予定向。把 $U_1\times U_2$ 上的向量丛
$E_1\oplus E_2$ 限制到 $U_1\times_BU_2$，其定向也随之诱导。
仓西邻域
\[
 (s_1\oplus s_2;E_1\oplus E_2\longrightarrow U_1\times_BU_2)
\]
通过
\[
 (E_1\oplus E_2)_*\times
 T_*(s_1\oplus s_2)^{-1}(0)=T_*(U_1\times_BU_2)
\]
定义 $(s_1\oplus s_2)^{-1}(0)$ 的定向。然后规定
\[
\begin{aligned}
 X_1\times_BX_2
 &=(-1)^{\operatorname{rank}E_2(\dim X_1-\dim B)}
   (s_1\oplus s_2)^{-1}(0)\\
 &=(-1)^{\operatorname{rank}E_2
   (\dim U_1-\operatorname{rank}E_1-\dim B)}
   (s_1\oplus s_2)^{-1}(0).
\end{aligned}
\]
这一约定不依赖于仓西邻域 $(s_i;E_i\to U_i)$ 的选择。
\end{enumerate}
\end{statement}

上述约定立即给出下列以后要用的基本公式。下文所谓
$(-1)$-定向同构（相应地，$(+1)$-定向同构），是指反定向同构
（相应地，保定向同构）。

\begin{lemma}{8.2.3}\label{lem:8-2-3}
\begin{enumerate}[label=(\arabic*),leftmargin=2.8em]
\item 假定 $X_1$、$X_2$ 有边界而 $\partial Y=\varnothing$。
对 $X_1\to Y$ 和 $X_2\to Y$，有
\[
 \partial(X_1\times_YX_2)
 =\partial X_1\times_YX_2
 \mathbin{\sqcup}(-1)^{x_1+y}(X_1\times_Y\partial X_2).
\]

\item \textup{（结合律）}
对 $X_1\to Y_1$、$X_2\to Y_1\times Y_2$ 和 $X_3\to Y_2$，有
\[
 (X_1\times_{Y_1}X_2)\times_{Y_2}X_3
 =X_1\times_{Y_1}(X_2\times_{Y_2}X_3).
\]

\item \textup{（迭代公式）}\liftIndex{I340}
对 $X_1\to Y_1\times Y_2$、$X_2\to Y_1$ 和 $X_3\to Y_2$，有
\[
 X_1\times_{Y_1\times Y_2}(X_2\times X_3)
 =(-1)^{y_2(y_1+x_2)}
 (X_1\times_{Y_1}X_2)\times_{Y_2}X_3.
\]

\item 设 $f_i\colon X_i\to X_i'$（$i=1,2$）是
$\epsilon(f_i)$-定向同构，而 $g\colon Y\to Y'$ 是
$\epsilon(g)$-定向同构。对 $X_i\to Y$ 和 $X_i'\to Y'$，诱导映射
\[
 f_1\times_gf_2\colon X_1\times_YX_2
 \longrightarrow X_1'\times_{Y'}X_2'
\]
是 $\epsilon(f_1)\epsilon(f_2)\epsilon(g)$-定向同构。
\end{enumerate}
\end{lemma}

\begin{proof}
除符号外各项都显然，故只核对符号。先考察所有映射都是流形之间浸没的
情形。

\textup{(1)} 按 \eqref{eq:8-2-1-5}--\eqref{eq:8-2-1-6}，令
\[
 X_1=X_1^\circ\times Y,
 \qquad X_2=Y\times{}^\circ X_2.
\]
再令
\[
 \partial X_1=(\partial X_1)^\circ\times Y,
 \qquad \partial X_2=Y\times{}^\circ(\partial X_2).
\]
由边界定向约定 \eqref{eq:8-2-1-1}，有
\[
\begin{aligned}
 \mathbb R_{X_1}\times(\partial X_1)^\circ\times Y
 &=X_1^\circ\times Y,\\
 \mathbb R_{X_2}\times\bigl(Y\times{}^\circ(\partial X_2)\bigr)
 &=Y\times{}^\circ X_2,
\end{aligned}
\]
其中 $\mathbb R_{X_i}$ 由外法向量定向。因此
\[
 \mathbb R_{X_1}\times(\partial X_1)^\circ=X_1^\circ,
 \qquad
 (-1)^y\bigl(\mathbb R_{X_2}\times{}^\circ(\partial X_2)\bigr)
 ={}^\circ X_2.
\]
所以
\[
\begin{aligned}
 &\mathbb R_{X_1\times_YX_2}\times\partial(X_1\times_YX_2)
 =X_1\times_YX_2=X_1^\circ\times Y\times{}^\circ X_2\\
 &\quad=
 \mathbb R_{X_1}\times(\partial X_1)^\circ\times Y\times{}^\circ X_2
 \mathbin{\sqcup}
 X_1^\circ\times Y\times
 (-1)^y\bigl(\mathbb R_{X_2}\times{}^\circ(\partial X_2)\bigr)\\
 &\quad=
 \mathbb R_{X_1}\times(\partial X_1\times_YX_2)
 \mathbin{\sqcup}
 (-1)^{x_1+y}\mathbb R_{X_2}
 \times(X_1\times_Y\partial X_2),
\end{aligned}
\]
这就证明了 (1)。

\textup{(2)} 令
\[
\begin{aligned}
 X_1&=X_1^\circ\times Y_1,\\
 X_2&=Y_1\times{}^\circ X_2
     =X_2^\circ\times Y_2
     =Y_1\times{}^\circ X_2^\circ\times Y_2,\\
 X_3&=Y_2\times{}^\circ X_3.
\end{aligned}
\]
于是 ${}^\circ X_2={}^\circ X_2^\circ\times Y_2$ 且
$X_2^\circ=Y_1\times{}^\circ X_2^\circ$。由于
\[
 X_1\times_{Y_1}X_2
 =X_1^\circ\times Y_1\times{}^\circ X_2
 =X_1^\circ\times Y_1\times{}^\circ X_2^\circ\times Y_2,
\]
我们有
\[
 (X_1\times_{Y_1}X_2)\times_{Y_2}X_3
 =X_1^\circ\times Y_1\times{}^\circ X_2^\circ
  \times Y_2\times{}^\circ X_3.
\]
另一方面，
\[
 X_2\times_{Y_2}X_3
 =X_2^\circ\times Y_2\times{}^\circ X_3
 =Y_1\times{}^\circ X_2^\circ\times Y_2\times{}^\circ X_3,
\]
故 (2) 的右端可写成
\[
 X_1\times_{Y_1}(X_2\times_{Y_2}X_3)
 =X_1^\circ\times Y_1\times{}^\circ X_2^\circ
  \times Y_2\times{}^\circ X_3.
\]
结合律随即成立。

\textup{(3)} 令
\[
 X_1=X_1^\circ\times Y_1\times Y_2,
 \qquad X_2=Y_1\times{}^\circ X_2,
 \qquad X_3=Y_2\times{}^\circ X_3.
\]
由于
\[
 X_2\times X_3
 =(-1)^{y_2(x_2-y_1)}
 Y_1\times Y_2\times{}^\circ X_2\times{}^\circ X_3,
\]
左端可写成
\begin{equation}
 X_1\times_{Y_1\times Y_2}(X_2\times X_3)
 =(-1)^{y_2(x_2-y_1)}
 X_1^\circ\times Y_1\times Y_2
 \times{}^\circ X_2\times{}^\circ X_3.
 \tag{8.2.4}\label{eq:8-2-4}
\end{equation}
另一方面，由于
\[
\begin{aligned}
 X_1\times_{Y_1}X_2
 &=(-1)^{y_1y_2}
 X_1^\circ\times Y_2\times Y_1\times{}^\circ X_2\\
 &=(-1)^{y_1y_2}(-1)^{y_2x_2}
 X_1^\circ\times Y_1\times{}^\circ X_2\times Y_2,
\end{aligned}
\]
右端可写成
\begin{align}
 &(X_1\times_{Y_1}X_2)\times_{Y_2}X_3\notag\\
 &\quad=(-1)^{y_2(y_1+x_2)}
 X_1^\circ\times Y_1\times{}^\circ X_2\times Y_2\times{}^\circ X_3
 \notag\\
 &\quad=(-1)^{y_2(y_1+x_2)}(-1)^{y_2(x_2-y_1)}
 X_1^\circ\times Y_1\times Y_2
 \times{}^\circ X_2\times{}^\circ X_3.
 \tag{8.2.5}\label{eq:8-2-5}
\end{align}
比较 \eqref{eq:8-2-4} 与 \eqref{eq:8-2-5} 即得公式。

\textup{(4)} 令
$X_1=X_1^\circ\times Y$、$X_2=Y\times{}^\circ X_2$，则
\[
 X_1\times_YX_2=X_1^\circ\times Y\times{}^\circ X_2.
\]
同样，
\[
 X_1'\times_{Y'}X_2'=X_1'{}^\circ\times Y'\times{}^\circ X_2'.
\]
由于 $Y'=\epsilon(g)Y$ 且 $X_i'=\epsilon(f_i)X_i$，有
\[
 X_1'{}^\circ=\epsilon(f_1)\epsilon(g)X_1^\circ,
 \qquad
 {}^\circ X_2'=\epsilon(f_2)\epsilon(g){}^\circ X_2.
\]
这就得到 (4)。

上述结论对带仓西结构空间同样成立。事实上，约定~\ref{conv:8-2-1}
的 (1)、(4) 正是为了让上述论证在作如下修改后仍然成立而设计的。
若 $\dim X<\dim Y$，就把 $X$ 看作带仓西结构空间
$(s;E\to U)$。当 $X$ 是 $Y$ 的子流形时，$X^\circ$ 在 $U=Y$ 上
扮演障碍丛 $E$ 的角色。

一般地，对带仓西结构空间，(1)--(4) 可证明如下。

\textup{(1)} 设 $X_i=(s_i;E_i\to U_i)$。由
约定~\ref{conv:8-2-1} 的 (1)、(4)，有
\[
\begin{aligned}
 \partial(X_1\times_YX_2)
 ={}&(-1)^{\delta_1}
 \bigl((s_1\oplus s_2)|_{\partial(U_1\times_YU_2)};\\
 &\qquad (E_1\oplus E_2)|_{\partial(U_1\times_YU_2)}
 \longrightarrow\partial(U_1\times_YU_2)\bigr),
\end{aligned}
\]
其中
\[
 \delta_1=\operatorname{rank}E_1+\operatorname{rank}E_2
 +\operatorname{rank}E_2(\dim X_1-\dim Y).
\]
同样，
\[
\begin{aligned}
 \partial X_1\times_YX_2
 ={}&(-1)^{\delta_2}
 \bigl((s_1\oplus s_2)|_{\partial U_1\times_YU_2};\\
 &\qquad (E_1\oplus E_2)|_{\partial U_1\times_YU_2}
 \longrightarrow\partial U_1\times_YU_2\bigr),
\end{aligned}
\]
其中
\[
 \delta_2=\operatorname{rank}E_2(\dim X_1-1-\dim Y)
 +\operatorname{rank}E_1,
\]
并且
\[
\begin{aligned}
 X_1\times_Y\partial X_2
 ={}&(-1)^{\delta_3}
 \bigl((s_1\oplus s_2)|_{U_1\times_Y\partial U_2};\\
 &\qquad (E_1\oplus E_2)|_{U_1\times_Y\partial U_2}
 \longrightarrow U_1\times_Y\partial U_2\bigr),
\end{aligned}
\]
其中
\[
 \delta_3=\operatorname{rank}E_2
 +\operatorname{rank}E_2(\dim X_1-\dim Y).
\]
对定向流形 $U_1,U_2$，有
\[
 \partial(U_1\times_YU_2)
 =\partial U_1\times_YU_2
 \mathbin{\sqcup}(-1)^{\delta_4}U_1\times_Y\partial U_2,
 \qquad \delta_4=\dim U_1+\dim Y.
\]
于是模 $2$ 有
\[
 \delta_1\equiv\delta_2,
 \qquad
 \delta_1+\delta_3+\delta_4\equiv\dim X_1+\dim Y.
\]
这里用到了
\[
 \dim U_1-\operatorname{rank}E_1=\dim X_1,
\]
即带仓西结构空间 $X_1=(s_1;E_1\to U_1)$ 的虚维数
（见定义 A1.5）。因此
\[
 \partial(X_1\times_YX_2)
 =\partial X_1\times_YX_2
 \mathbin{\sqcup}(-1)^{\dim X_1+\dim Y}
 X_1\times_Y\partial X_2.
\]

\textup{(2)} 由约定~\ref{conv:8-2-1} (4)，有
\[
\begin{aligned}
 &(X_1\times_{Y_1}X_2)\times_{Y_2}X_3\\
 &\quad=(-1)^{\delta_5}
 (s_1\oplus s_2;E_1\oplus E_2\to U_1\times_{Y_1}U_2)
 \times_{Y_2}X_3\\
 &\quad=(-1)^{\delta_5+\delta_6}
 \bigl(s_1\oplus s_2\oplus s_3;
 E_1\oplus E_2\oplus E_3
 \to(U_1\times_{Y_1}U_2)\times_{Y_2}U_3\bigr),
\end{aligned}
\]
其中
\[
\begin{aligned}
 \delta_5&=\operatorname{rank}E_2(\dim X_1-\dim Y_1),\\
 \delta_6&=\operatorname{rank}E_3
 (\dim X_1+\dim X_2-\dim Y_1-\dim Y_2).
\end{aligned}
\]
同样，
\[
\begin{aligned}
 &X_1\times_{Y_1}(X_2\times_{Y_2}X_3)\\
 &\quad=(-1)^{\delta_7}X_1\times_{Y_1}
 (s_2\oplus s_3;E_2\oplus E_3\to U_2\times_{Y_2}U_3)\\
 &\quad=(-1)^{\delta_7+\delta_8}
 \bigl(s_1\oplus s_2\oplus s_3;
 E_1\oplus E_2\oplus E_3
 \to U_1\times_{Y_1}(U_2\times_{Y_2}U_3)\bigr),
\end{aligned}
\]
其中
\[
\begin{aligned}
 \delta_7&=\operatorname{rank}E_3(\dim X_2-\dim Y_2),\\
 \delta_8&=(\operatorname{rank}E_2+\operatorname{rank}E_3)
 (\dim X_1-\dim Y_1).
\end{aligned}
\]
对定向流形 $U_1,U_2,U_3$，上面已经证明
\[
 (U_1\times_{Y_1}U_2)\times_{Y_2}U_3
 =U_1\times_{Y_1}(U_2\times_{Y_2}U_3).
\]
又因 $\delta_5+\delta_6=\delta_7+\delta_8$，故
\[
 (X_1\times_{Y_1}X_2)\times_{Y_2}X_3
 =X_1\times_{Y_1}(X_2\times_{Y_2}X_3).
\]

\textup{(3)} 由约定~\ref{conv:8-2-1} (4)，有
\[
\begin{aligned}
 &X_1\times_{Y_1\times Y_2}(X_2\times X_3)\\
 &\quad=(-1)^{\delta_9}X_1\times_{Y_1\times Y_2}
 (s_2\oplus s_3;E_2\oplus E_3\to U_2\times U_3)\\
 &\quad=(-1)^{\delta_9+\delta_{10}}
 \bigl(s_1\oplus s_2\oplus s_3;
 E_1\oplus E_2\oplus E_3
 \to U_1\times_{Y_1\times Y_2}(U_2\times U_3)\bigr),
\end{aligned}
\]
其中
\[
\begin{aligned}
 \delta_9&=\operatorname{rank}E_3\cdot\dim X_2,\\
 \delta_{10}&=(\operatorname{rank}E_2+\operatorname{rank}E_3)
 (\dim X_1-\dim Y_1-\dim Y_2).
\end{aligned}
\]
另一方面，
\[
\begin{aligned}
 &(X_1\times_{Y_1}X_2)\times_{Y_2}X_3\\
 &\quad=(-1)^{\delta_{11}}
 (s_1\oplus s_2;E_1\oplus E_2\to U_1\times_{Y_1}U_2)
 \times_{Y_2}X_3\\
 &\quad=(-1)^{\delta_{11}+\delta_{12}}
 \bigl(s_1\oplus s_2\oplus s_3;
 E_1\oplus E_2\oplus E_3
 \to(U_1\times_{Y_1}U_2)\times_{Y_2}U_3\bigr),
\end{aligned}
\]
其中
\[
\begin{aligned}
 \delta_{11}&=\operatorname{rank}E_2(\dim X_1-\dim Y_1),\\
 \delta_{12}&=\operatorname{rank}E_3
 (\dim X_1+\dim X_2-\dim Y_1-\dim Y_2).
\end{aligned}
\]
回忆对定向流形 $U_1,U_2,U_3$，已经证明
\[
 U_1\times_{Y_1\times Y_2}(U_2\times U_3)
 =(-1)^{\delta_{13}}(U_1\times_{Y_1}U_2)\times_{Y_2}U_3,
\]
其中
\[
 \delta_{13}=\dim Y_2(\dim Y_1+\dim U_2).
\]
由于模 $2$ 有
\[
 \delta_9+\delta_{10}+\delta_{11}+\delta_{12}+\delta_{13}
 \equiv\dim Y_2(\dim Y_1+\dim X_2),
\]
我们得到
\[
 X_1\times_{Y_1\times Y_2}(X_2\times X_3)
 =(-1)^{\dim Y_2(\dim Y_1+\dim X_2)}
 (X_1\times_{Y_1}X_2)\times_{Y_2}X_3.
\]

\textup{(4)} 约定~\ref{conv:8-2-1} (4) 对带仓西结构空间的纤维积
$X_1\times_YX_2$ 所定义的定向，只依赖于 $X_1,X_2$ 的仓西结构定向
以及 $Y$ 的定向；换言之，它不依赖于带定向局部坐标图的选择。因此
(4) 显然成立。
\end{proof}

\begin{remark}{8.2.6}\label{rem:8-2-6}
利用上述纤维积的写法，可以把下面的纤维积定向简明地写出。设
\[
 f_k\colon X\longrightarrow Y_k,
 \qquad g_k\colon X_k\longrightarrow Y_k
 \qquad(k=1,\ldots,\ell).
\]
对 $k=1,\ldots,\ell$，令
\[
 X=:X^\circ\times\prod_{k=1}^{\ell}Y_k,
 \qquad
 X_k=:Y_k\times{}^\circ X_k.
\]
那么纤维积的定向由下式给出：
\begin{equation}
 X_{(f_1,\ldots,f_\ell)}\!\times_{g_1\times\cdots\times g_\ell}
 \left(\prod_{k=1}^{\ell}X_k\right)
 =X^\circ\times\prod_{k=1}^{\ell}(Y_k\times{}^\circ X_k).
 \tag{8.2.7}\label{eq:8-2-7}
\end{equation}

为证明 \eqref{eq:8-2-7}，令
\[
 \prod_{k=1}^{\ell}X_k
 =:\left(\prod_{k=1}^{\ell}Y_k\right)
 \times{}^\circ\left(\prod_{k=1}^{\ell}X_k\right).
\]
由纤维积定向的定义，
\[
 X_{(f_1,\ldots,f_\ell)}\!\times_{g_1\times\cdots\times g_\ell}
 \left(\prod_{k=1}^{\ell}X_k\right)
 =X^\circ\times\left(\prod_{k=1}^{\ell}Y_k\right)
 \times{}^\circ\left(\prod_{k=1}^{\ell}X_k\right).
\]
另一方面，容易看出
\[
 {}^\circ\left(\prod_{k=1}^{\ell}X_k\right)
 =(-1)^{\delta_1}\prod_{k=1}^{\ell}{}^\circ X_k,
 \qquad
 \delta_1=\sum_{k=2}^{\ell}y_k
 \sum_{j=1}^{k-1}(x_j-y_j).
\]
因而
\[
\begin{aligned}
 &X_{(f_1,\ldots,f_\ell)}\!\times_{g_1\times\cdots\times g_\ell}
 \left(\prod_{k=1}^{\ell}X_k\right)\\
 &\quad=(-1)^{\delta_1}X^\circ
 \times\left(\prod_{k=1}^{\ell}Y_k\right)
 \times\left(\prod_{k=1}^{\ell}{}^\circ X_k\right)\\
 &\quad=(-1)^{\delta_1+\delta_2}X^\circ
 \times\prod_{k=1}^{\ell}(Y_k\times{}^\circ X_k).
\end{aligned}
\]
显然 $\delta_1=\delta_2$，故得到公式 \eqref{eq:8-2-7}。
\end{remark}
